\documentclass[11pt]{article}
\usepackage{shared/luxcover}
\usepackage{amsmath, amssymb, algpseudocode, graphicx, hyperref, booktabs}

\title{Signed Software, Reproducible Builds, and Quantum-Resilient Release Pipelines}
\author{
Zach Kelling\thanks{Corresponding author: zach@lux.network}\\
\textit{Lux Industries, Inc.}\\
\texttt{research@lux.network}
}
\date{2025}

\usepackage{amsthm}
\newtheorem{definition}{Definition}
\newtheorem{theorem}{Theorem}
\begin{document}
\luxcoverpage

\begin{abstract}
A blockchain node binary is the ultimate trust anchor: if the binary is compromised, no amount of cryptographic protocol design matters. We describe Lux Network's release pipeline, which enforces deterministic builds, ML-DSA (FIPS 204) signed binaries, and an attestation chain from source commit to deployed container image. Every release artifact carries a hybrid BLS+ML-DSA signature over a content-addressed manifest that includes source commit hash, build environment hash, compiler version, and dependency lockfile digest. Reproducibility is verified by independent rebuilders, and attestations are published to a transparency log anchored in Q-Chain. The pipeline covers \texttt{luxd}, EVM plugins, and subnet VM binaries across Linux (amd64, arm64) targets.
\end{abstract}

\section{Introduction}

Supply chain attacks on software dependencies have escalated from theoretical concern to operational reality. The SolarWinds incident (2020), the xz/liblzma backdoor (2024), and numerous npm/PyPI compromises demonstrate that build pipelines are high-value targets. For blockchain infrastructure, a compromised node binary can steal funds, halt consensus, or silently corrupt state.

Lux Network addresses this with three interlocking guarantees:
\begin{enumerate}
\item \textbf{Deterministic builds}: Given the same source commit and build environment, any party produces the identical binary.
\item \textbf{Quantum-resilient signatures}: Release artifacts are signed with ML-DSA-65 (FIPS 204) in addition to classical Ed25519, ensuring signature validity even against future quantum adversaries.
\item \textbf{Attestation transparency}: Every step from source to deployment is recorded in an append-only log anchored in Q-Chain, enabling public audit.
\end{enumerate}

\subsection{Scope}

The pipeline covers:
\begin{itemize}
\item \texttt{luxd} node binary (Go, \texttt{CGO\_ENABLED=0})
\item EVM plugin binaries (Go, compiled against \texttt{luxfi/evm})
\item Subnet VM binaries (application-specific)
\item Container images (\texttt{ghcr.io/luxfi/*}, \texttt{--platform linux/amd64})
\end{itemize}

\section{Deterministic Build System}

\subsection{Build Environment Specification}

Each release specifies a locked build environment:

\begin{verbatim}
build-env.lock:
  go_version: 1.22.4
  os: linux
  arch: amd64
  cgo: disabled
  ldflags: -s -w -X main.version=v1.23.23
  go_sum_hash: sha256:<digest>
  toolchain_hash: sha256:<digest>
\end{verbatim}

The build environment itself is a content-addressed container image. The image is built once, signed, and reused for all builds at that Go version. Changing the Go version requires a new environment image with a new signature.

\subsection{Reproducibility Protocol}

\begin{algorithmic}[1]
\State \textbf{function} \textsc{Build}(commit, envImage)
\State $source \gets \text{GitCheckout}(commit)$
\State $deps \gets \text{GoModDownload}(source)$
\State $binary \gets \text{GoBuild}(source, deps, envImage)$
\State $digest \gets \text{SHA256}(binary)$
\State \Return $(binary, digest)$
\State \textbf{end function}
\end{algorithmic}

Two properties ensure reproducibility:
\begin{enumerate}
\item \textbf{No embedded timestamps}: Build flags strip debug info (\texttt{-s -w}) and inject only the version string and commit hash. No build time, hostname, or path is embedded.
\item \textbf{Locked dependencies}: \texttt{go.sum} is committed and verified. The dependency fetch is hermetic---no network access during compilation.
\end{enumerate}

\subsection{Independent Rebuilder Verification}

At least two independent rebuilders (community-operated, not Lux Industries) reproduce each release. A release is considered verified when $\geq 2$ independent rebuilders produce the identical binary digest.

\begin{table}[h]
\centering
\begin{tabular}{@{}llr@{}}
\toprule
\textbf{Artifact} & \textbf{Target} & \textbf{Binary Size} \\
\midrule
luxd & linux/amd64 & $\sim$48 MB \\
luxd & linux/arm64 & $\sim$45 MB \\
evm-plugin & linux/amd64 & $\sim$32 MB \\
\bottomrule
\end{tabular}
\caption{Typical artifact sizes for Lux v1.23.x releases.}
\label{tab:artifacts}
\end{table}

\section{Hybrid Signature Scheme}

\subsection{Release Manifest}

Each release produces a signed manifest:

\begin{verbatim}
manifest.json:
{
  "version": "v1.23.23",
  "commit": "abc123def456",
  "artifacts": [
    {
      "name": "luxd-linux-amd64",
      "sha256": "<digest>",
      "size": 50331648
    }
  ],
  "build_env_hash": "<digest>",
  "go_version": "1.22.4",
  "go_sum_hash": "<digest>",
  "timestamp": "2025-03-01T00:00:00Z"
}
\end{verbatim}

\subsection{Dual Signature}

The manifest is signed with both Ed25519 (classical) and ML-DSA-65 (post-quantum):

\begin{align}
\sigma_{\text{ed25519}} &= \text{Sign}_{\text{Ed25519}}(sk_c, H(\text{manifest})) \\
\sigma_{\text{mldsa}} &= \text{Sign}_{\text{MLDSA}}(sk_{pq}, H(\text{manifest})) \\
\text{release} &= (\text{manifest}, \sigma_{\text{ed25519}}, \sigma_{\text{mldsa}})
\end{align}

The signing keys are held in an HSM \cite{hsm-boundary}. The ML-DSA key is generated and stored in the same HSM boundary as the Ed25519 key. Neither key is extractable.

\subsection{Verification}

Validators verify release signatures before applying upgrades:

\begin{algorithmic}[1]
\State \textbf{function} \textsc{VerifyRelease}(release)
\State $h \gets \text{SHA256}(release.\text{manifest})$
\State $v_c \gets \text{VerifyEd25519}(pk_c, \sigma_{\text{ed25519}}, h)$
\State $v_{pq} \gets \text{VerifyMLDSA}(pk_{pq}, \sigma_{\text{mldsa}}, h)$
\State \Return $v_c \land v_{pq}$
\State \textbf{end function}
\end{algorithmic}

\section{Attestation Transparency}

\subsection{Attestation Chain}

Every build step produces an attestation record:

\begin{table}[h]
\centering
\small
\begin{tabular}{@{}lll@{}}
\toprule
\textbf{Step} & \textbf{Attestation} & \textbf{Signer} \\
\midrule
Source & Commit signed (GPG/SSH) & Developer \\
Build & Binary digest + env hash & CI runner \\
Rebuild & Matching digest confirmation & Independent rebuilder \\
Sign & Manifest + dual signature & Release HSM \\
Publish & Container image digest & Registry (GHCR) \\
Deploy & Running image digest & Validator node \\
\bottomrule
\end{tabular}
\caption{Attestation chain from source to deployment.}
\label{tab:attestation}
\end{table}

\subsection{Q-Chain Anchor}

Attestation records are batched and their Merkle root is submitted to Q-Chain as a governance transaction. This provides a tamper-evident, publicly auditable record that cannot be rewritten without breaking Q-Chain consensus.

Validators can verify that the binary they are running corresponds to a release whose attestation root is recorded on-chain:

\begin{algorithmic}[1]
\State \textbf{function} \textsc{VerifyDeployment}(binaryDigest)
\State $root \gets \text{QChain.LatestAttestationRoot}()$
\State $proof \gets \text{FetchMerkleProof}(binaryDigest, root)$
\State \Return $\text{VerifyMerkleProof}(proof, root, binaryDigest)$
\State \textbf{end function}
\end{algorithmic}

\section{Container Image Pipeline}

Container images are built from the signed binary, not from source. The Dockerfile is minimal:

\begin{verbatim}
FROM debian:bookworm-slim
COPY --from=signed luxd /usr/local/bin/luxd
COPY --from=signed plugins/ /usr/local/lib/luxd/plugins/
ENTRYPOINT ["luxd"]
\end{verbatim}

Images are built using Kaniko (no Docker daemon) inside the Kubernetes cluster, pushed to \texttt{ghcr.io/luxfi/node:<tag>}, and signed with Cosign using the same ML-DSA key. The image digest is included in the attestation chain.

Images are built for \texttt{linux/amd64} only. Multi-arch images are produced by CI runners on the target architecture, not by cross-compilation or emulation.

\section{Threat Model and Mitigations}

\begin{table}[h]
\centering
\small
\begin{tabular}{@{}lll@{}}
\toprule
\textbf{Threat} & \textbf{Mitigation} & \textbf{Detection} \\
\midrule
Compromised CI runner & Reproducible builds & Rebuilder mismatch \\
Dependency backdoor & Locked go.sum, hermetic build & Hash mismatch on rebuild \\
Signing key theft & HSM-bound keys, dual signature & Requires both key types \\
Quantum signature forge & ML-DSA alongside Ed25519 & Hybrid check fails \\
Binary swap post-sign & Content-addressed manifest & Digest verification \\
\bottomrule
\end{tabular}
\caption{Threat mitigations in the release pipeline.}
\label{tab:threats}
\end{table}

\section{Conclusion}

The distance from source code to running validator is the most security-critical path in a blockchain system. Lux's release pipeline closes every link in this chain: deterministic builds eliminate the ``it built differently'' class of attacks, hybrid ML-DSA+Ed25519 signatures protect against both current and future adversaries, and Q-Chain-anchored attestation transparency makes every step publicly auditable. Independent rebuilders provide the final check: if the binary cannot be reproduced, it is not released.

\begin{thebibliography}{99}

\bibitem{hsm-boundary}
Kelling, Z. (2025).
\textit{Secure Enclave, HSM, and Threshold Boundary Design}.
Lux Industries Research.

\bibitem{quasar}
Kelling, Z. (2025).
\textit{Quasar: Quantum-Secure Multi-Engine Consensus with Triple-Proof Quantum Finality}.
Lux Industries Research.

\bibitem{nist-fips204}
NIST (2024).
\textit{FIPS 204: Module-Lattice-Based Digital Signature Standard (ML-DSA)}.
National Institute of Standards and Technology.

\bibitem{reproducible-builds}
Reproducible Builds Project (2024).
\textit{Reproducible Builds: A Set of Software Development Practices}.
\url{https://reproducible-builds.org/}.

\bibitem{sigstore}
Sigstore Project (2024).
\textit{Sigstore: A New Standard for Signing, Verifying, and Protecting Software}.
Linux Foundation.

\bibitem{slsa}
SLSA Contributors (2024).
\textit{Supply-chain Levels for Software Artifacts (SLSA)}.
\url{https://slsa.dev/}.

\end{thebibliography}

\end{document}
